% =============================================================================
% 课程论文LaTeX模板
% FormalMath项目学术写作支持系统
% 适用：数学课程论文、学期论文（5-20页）
% =============================================================================

\documentclass[12pt,a4paper]{article}

% ==================== 宏包引入 ====================
% 中文支持（使用XeLaTeX编译）
\usepackage{ctex}

% 页面设置
\usepackage{geometry}
\geometry{left=2.5cm,right=2.5cm,top=2.5cm,bottom=2.5cm}

% 数学宏包
\usepackage{amsmath,amssymb,amsthm}
\usepackage{mathtools}

% 定理环境设置
\theoremstyle{definition}
\newtheorem{theorem}{定理}[section]
\newtheorem{lemma}[theorem]{引理}
\newtheorem{proposition}[theorem]{命题}
\newtheorem{corollary}[theorem]{推论}
\newtheorem{definition}[theorem]{定义}
\newtheorem{example}[theorem]{例子}
\newtheorem{remark}[theorem]{注记}

% 图表支持
\usepackage{graphicx}
\usepackage{tikz}
\usepackage{tikz-cd}
\usetikzlibrary{arrows,shapes,positioning}

% 参考文献
\usepackage[numbers]{natbib}
\bibliographystyle{amsplain}

% 超链接
\usepackage{hyperref}
\hypersetup{
    colorlinks=true,
    linkcolor=blue,
    citecolor=blue,
    urlcolor=blue
}

% 页眉页脚
\usepackage{fancyhdr}
\pagestyle{fancy}
\fancyhf{}
\fancyhead[L]{\small 课程名称}
\fancyhead[R]{\small 论文标题}
\fancyfoot[C]{\thepage}
\renewcommand{\headrulewidth}{0.4pt}

% 行间距
\usepackage{setspace}
\onehalfspacing

% 列表设置
\usepackage{enumitem}

% ==================== 文档信息 ====================
\title{\textbf{论文标题}}
\author{作者姓名 \\ 学号：12345678 \\ 学校/院系}
\date{\today}

% ==================== 正文开始 ====================
\begin{document}

\maketitle

% ==================== 摘要 ====================
\begin{abstract}
这是摘要。摘要应该独立概述全文，包括研究背景、主要内容、方法和主要结果。
长度通常在150-250词之间。

\vspace{1em}
\noindent\textbf{关键词：} 关键词1；关键词2；关键词3
\end{abstract}

% ==================== 正文 ====================
\section{引言}

引言应该：
\begin{enumerate}
    \item 引入研究主题，说明背景和动机
    \item 陈述本文要解决的问题
    \item 简要回顾相关文献
    \item 概述论文的结构
\end{enumerate}

\subsection{背景}

从读者熟悉的内容开始，逐步引入主题。

\subsection{问题陈述}

明确说明本文要研究的具体问题。

\subsection{论文结构}

本文结构如下。第\ref{sec:prelim}节介绍必要的预备知识。第\ref{sec:main}节呈现主要结果。第\ref{sec:applications}节讨论应用和例子。最后，第\ref{sec:conclusion}节总结全文。

% ==================== 预备知识 ====================
\section{预备知识}\label{sec:prelim}

本节回顾必要的定义和已知结果。

\begin{definition}\label{def:example}
这是一个定义示例。
\end{definition}

\begin{theorem}[已知定理]\label{thm:known}
这是一个已知的定理陈述。
\end{theorem}

\begin{proof}
这里可以给出证明或引用参考文献。
\end{proof}

% ==================== 主要结果 ====================
\section{主要结果}\label{sec:main}

本节陈述并证明本文的主要结果。

\begin{theorem}[主要定理]\label{thm:main}
这是本文的主要定理。
\end{theorem}

\begin{proof}
证明应该清晰、完整，每一步都有明确说明。
\end{proof}

\begin{corollary}\label{cor:main}
这是从主要定理得出的推论。
\end{corollary}

% ==================== 应用与例子 ====================
\section{应用与例子}\label{sec:applications}

本节通过具体例子说明理论的应用。

\begin{example}
这是一个具体的例子。
\end{example}

\begin{figure}[htbp]
    \centering
    \begin{tikzpicture}
        % 在这里添加图形代码
        \draw[->] (-2,0) -- (2,0) node[right] {$x$};
        \draw[->] (0,-1) -- (0,2) node[above] {$y$};
        \draw[domain=-1.5:1.5,smooth,variable=\x,blue] plot (\x,{\x*\x});
    \end{tikzpicture}
    \caption{函数 $y=x^2$ 的图像}
    \label{fig:example}
\end{figure}

% ==================== 结论 ====================
\section{结论}\label{sec:conclusion}

结论应该：
\begin{itemize}
    \item 总结主要结果
    \item 讨论局限性和未解决问题
    \item 展望未来研究方向
\end{itemize}

% ==================== 参考文献 ====================
\newpage
\bibliography{references}

% ==================== 附录（可选） ====================
\appendix
\section{补充证明}

这里可以放置详细的补充材料。

\end{document}
% =============================================================================
% 使用说明：
% 1. 确保已安装TeX发行版（TeX Live或MiKTeX）
% 2. 使用XeLaTeX编译
% 3. 如需参考文献，创建references.bib文件
% 4. 根据具体课程要求调整页眉信息
% =============================================================================
